\documentclass[11pt]{article}

\usepackage[margin=1in]{geometry}
\usepackage{amsmath, amssymb, amsthm, mathtools}
\usepackage{hyperref}
\usepackage{enumitem}

\title{A Correctness Note for the Factor Library Reward Mechanism}
\author{}
\date{April 15, 2026}

\newtheorem{theorem}{Theorem}
\newtheorem{proposition}[theorem]{Proposition}
\newtheorem{lemma}[theorem]{Lemma}
\newtheorem{corollary}[theorem]{Corollary}
\theoremstyle{remark}
\newtheorem{remark}[theorem]{Remark}

\newcommand{\Fp}{\mathbb{F}_p}
\newcommand{\Z}{\mathbb{Z}}
\newcommand{\B}{\mathcal{B}}
\newcommand{\Llib}{\mathcal{L}}
\newcommand{\Gset}{\mathcal{G}}
\newcommand{\States}{\mathcal{S}}

\begin{document}
\maketitle

\begin{abstract}
This note formalizes the factor-library mechanism implemented in
\texttt{src/environment/factor\_library.py} and
\texttt{src/environment/circuit\_game.py}. We prove five implementation-faithful
claims. First, the factor library never changes the transition rule or the exact
success condition of the circuit-construction environment; it only changes
reward shaping and cross-episode bookkeeping. Second, every initial factor
subgoal returned by the library is a genuine nontrivial divisor of the target
polynomial modulo $p$. Third, every dynamically discovered residual or exact
quotient subgoal is algebraically valid. Fourth, both completion bonuses are
exact certificates that the current state is one operation away from the target:
the additive bonus certifies one addition away, and the multiplicative bonus
certifies one multiplication away. Fifth, every polynomial stored in the session
library has a witnessed circuit construction from some successful episode.
\end{abstract}

\section{Setup}

Let $p$ be the environment modulus, let $n$ be the number of variables, and let
$d$ be the configured per-variable degree cap. The implementation stores each
polynomial as a dense coefficient tensor of shape $(d+1)^n$, so the exact
algebraic object manipulated by \texttt{FastPoly} is the truncated ring
\[
R_{p,d,n}
:=
\Fp[x_0,\dots,x_{n-1}]
\big/
\langle x_0^{d+1},\dots,x_{n-1}^{d+1}\rangle.
\]
Addition is exact in this quotient ring, and multiplication is exact in this
quotient ring because the code performs convolution and then discards monomials
whose exponent in some variable exceeds $d$.

Let
\[
\B := \{x_0,\dots,x_{n-1},1\} \subset R_{p,d,n}
\]
be the base nodes present at the start of every episode. A state is an ordered
tuple
\[
s_t = (v_0,\dots,v_{m_t-1}) \in \States,
\]
where the first $n+1$ entries are the base nodes and each later node is formed
from earlier nodes by one binary operation $+$ or $\times$.

An action is a triple $(\star,i,j)$ with $\star \in \{+,\times\}$ and
$0 \le i \le j < m_t$. The next node is
\[
v_{m_t} = v_i \star v_j \in R_{p,d,n}.
\]
For a fixed target polynomial $T \in R_{p,d,n}$, the success predicate is
\[
\mathrm{Succ}(s_t) \iff \exists k < m_t \text{ such that } v_k = T.
\]

The session-level factor library is a map
\[
\Llib : R_{p,d,n} \setminus \B \to \mathbb{N}
\]
that records polynomials seen in successful episodes, together with the first
step index at which they were observed in that episode. During a single episode,
the environment also maintains a subgoal set $\Gset_t \subseteq R_{p,d,n}$ and
a set of already-claimed subgoals to prevent double counting.

\paragraph{Reduction maps.}
Let
\[
\rho_p : \Z[x_0,\dots,x_{n-1}] \to \Fp[x_0,\dots,x_{n-1}]
\]
be coefficient-wise reduction modulo $p$, and let
\[
\tau_d : \Fp[x_0,\dots,x_{n-1}] \to R_{p,d,n}
\]
be the quotient map. Their composition $\pi := \tau_d \circ \rho_p$ is exactly
what the code implements when it converts a SymPy polynomial into a
\texttt{FastPoly}.

\section{Reward Decomposition}

Ignoring action masking and termination caps, the reward implemented in
\texttt{CircuitGame.step()} can be written as
\begin{align*}
r_t
&=
c_{\mathrm{step}}
+ \mathbf{1}\{\mathrm{Succ}(s_{t+1})\} \, c_{\mathrm{succ}} \\
&\quad
+ \mathbf{1}\{\neg \mathrm{Succ}(s_{t+1})\}
\bigl(\gamma \Phi(s_{t+1}) - \Phi(s_t)\bigr) \\
&\quad
+ c_{\mathrm{fac}} I^{\mathrm{fac}}_t
+ c_{\mathrm{lib}} I^{\mathrm{lib}}_t
+ c_{\mathrm{comp}} I^{+}_t
+ c_{\mathrm{comp}} I^{\times}_t,
\end{align*}
where:
\begin{itemize}[leftmargin=2em]
\item $\Phi$ is the term-similarity potential used by the existing shaping rule;
\item $I^{\mathrm{fac}}_t$ indicates that the newly built node belongs to the
current episode's active subgoal set and has not been rewarded before;
\item $I^{\mathrm{lib}}_t$ indicates that this subgoal was also in the
cross-episode library;
\item $I_t^+$ and $I_t^\times$ are the additive and multiplicative completion
indicators.
\end{itemize}

The potential-based term is standard shaping in the sense of
Ng--Harada--Russell (1999). The factor-library terms are different: they are
not policy-invariance corrections, but explicit guidance toward reusable
intermediates. The goal of this note is therefore \emph{soundness}, not
policy-invariance: we prove that these bonuses point toward algebraically valid
subcomputations and do not alter the exact success criterion.

\section{State-Semantics Invariants}

\begin{lemma}[Circuit semantics]
Every node $v_k$ created by the environment equals the polynomial computed by
its corresponding arithmetic circuit in $R_{p,d,n}$.
\end{lemma}

\begin{proof}
The claim is immediate by induction on the creation order. The base nodes are
exactly $x_0,\dots,x_{n-1},1$. If the claim holds for all previously created
nodes and the environment applies $(\star,i,j)$, then the new node is defined
to be $v_i \star v_j$ using the ring operations implemented by
\texttt{FastPoly}. Those operations are the quotient-ring operations of
$R_{p,d,n}$, so the semantic equation is preserved.
\end{proof}

\begin{proposition}[The factor library does not change what counts as success]
Fix an episode target $T$ and an action sequence
$a_0,\dots,a_{K-1}$. Suppose the environment is run twice from the same initial
state and with the same action sequence, once with the factor-library logic
disabled and once with it enabled. Then both runs produce the same sequence of
nodes and the same terminal success/failure outcome. The only differences are
reward values and auxiliary bookkeeping sets.
\end{proposition}

\begin{proof}
In \texttt{CircuitGame.step()}, node creation depends only on the decoded action,
the current node list, and the deterministic \texttt{FastPoly} operations. The
success predicate is the equality test $v_{\mathrm{new}} = T$, and the done flag
depends only on success, the step cap, and the node cap. The factor-library
logic is executed only after the new node has been created and success has been
checked. It modifies reward and updates auxiliary sets
$(\Gset_t,\Llib)$, but it does not rewrite the node list, the target, or the
success predicate.
\end{proof}

\section{Soundness of Initial Factor Subgoals}

At reset time, the code converts the target $T$ into a representative
$\widetilde{T} \in \Z[x_0,\dots,x_{n-1}]$ with coefficients in
$\{0,1,\dots,p-1\}$ and asks SymPy for a factorization
\[
\widetilde{T}
=
c \prod_{\ell=1}^m \widetilde{f}_{\ell}^{e_\ell}
\qquad
\text{in }
\Z[x_0,\dots,x_{n-1}].
\]
Each candidate factor is then reduced by $\pi$, and the implementation discards
scalars, base nodes, zero polynomials, duplicates, and the target itself.

\begin{proposition}[Every returned initial factor is a true divisor mod $p$]
Let $g \in R_{p,d,n}$ be any element returned by
\texttt{factorize\_target(T)}. Then there exists
$h \in R_{p,d,n}$ such that
\[
T = g h
\qquad \text{in } R_{p,d,n}.
\]
Moreover, $g$ is nonzero, non-scalar, not a base node, and not equal to $T$.
\end{proposition}

\begin{proof}
By construction, $g = \pi(\widetilde{f}_i)$ for some integer factor
$\widetilde{f}_i$ in the above factorization, and the filters guarantee that
$g$ is nonzero, non-scalar, not in $\B$, and not equal to $T$. Let
\[
\widetilde{h}
:=
c \, \widetilde{f}_i^{e_i-1} \prod_{\ell \neq i} \widetilde{f}_{\ell}^{e_\ell}
\in \Z[x_0,\dots,x_{n-1}].
\]
Then $\widetilde{T} = \widetilde{f}_i \widetilde{h}$ in the integer polynomial
ring. Applying $\pi$ gives
\[
T
=
\pi(\widetilde{T})
=
\pi(\widetilde{f}_i)\pi(\widetilde{h})
=
g h
\]
with $h := \pi(\widetilde{h}) \in R_{p,d,n}$.
\end{proof}

\begin{remark}
Factoring over $\Z$ and then reducing modulo $p$ is therefore sound for
discovering divisors: each individual factor still divides the target after
reduction. The choice of factoring over $\Z$ is about numerical robustness and
implementation convenience, not about changing the divisor relation.
\end{remark}

\section{Soundness of Dynamic Discovery}

Dynamic discovery is gated by the condition that the newly built node
$v := v_{\mathrm{new}}$ is already in the session library. This gating is purely
computational; it reduces SymPy calls. The algebraic claims below do not depend
on why the gate was opened.

\subsection{Residual subgoals}

The environment computes the residual
\[
R_v := T - v \in R_{p,d,n}.
\]
It may add $R_v$ itself to the active subgoal set (unless it is zero or a base
node), and it may also factorize $R_v$ to add nontrivial factors of $R_v$.

\begin{proposition}[Residual soundness]
If the environment adds the residual $R_v$ to the active subgoal set, then
\[
v + R_v = T
\qquad \text{in } R_{p,d,n}.
\]
If the environment additionally adds a factor $g$ returned by
\texttt{factorize\_poly(R\_v)}, then $g$ is a nontrivial divisor of $R_v$ in
$R_{p,d,n}$.
\end{proposition}

\begin{proof}
The identity $v + R_v = T$ is true by definition of $R_v$. The second claim is
exactly the previous proposition applied to $R_v$ instead of $T$, because
\texttt{factorize\_poly} uses the same factorization-and-filtering procedure on
its input polynomial.
\end{proof}

\subsection{Exact quotients}

The exact-quotient routine converts $T$ and $v$ into integer representatives
$\widetilde{T},\widetilde{v}$, computes
\[
\widetilde{T} = \widetilde{v}\widetilde{q} + \widetilde{r}
\qquad \text{in } \Z[x_0,\dots,x_{n-1}],
\]
and accepts the quotient precisely when $\pi(\widetilde{r}) = 0$.

\begin{proposition}[Exact quotient soundness]
If \texttt{exact\_quotient(T,v)} returns $q \in R_{p,d,n}$, then
\[
T = v q
\qquad \text{in } R_{p,d,n}.
\]
If the environment further factorizes $q$ and adds some factor $g$, then $g$ is
a nontrivial divisor of $q$ in $R_{p,d,n}$.
\end{proposition}

\begin{proof}
By the division identity over $\Z$,
\[
\widetilde{T} = \widetilde{v}\widetilde{q} + \widetilde{r}.
\]
Applying $\pi$ and using the acceptance condition $\pi(\widetilde{r}) = 0$
yields
\[
T = v q + 0 = v q
\qquad \text{in } R_{p,d,n},
\]
where $q = \pi(\widetilde{q})$. The factor statement again follows from the
soundness of \texttt{factorize\_poly} applied to the polynomial $q$.
\end{proof}

\section{Completion Bonuses}

\subsection{Additive completion}

The additive completion bonus fires when the residual $R_v := T-v$ is already
present among the old nodes of the current state.

\begin{theorem}[Additive completion is an exact one-step certificate]
Suppose the additive completion bonus fires at a state whose newly created node
is $v$. Then there exists an already existing node $u$ such that
\[
u = T-v
\qquad \text{and hence} \qquad
u + v = T
\]
in $R_{p,d,n}$. Therefore the current state is exactly one addition away from
the target.
\end{theorem}

\begin{proof}
By the firing condition, there is an existing node $u$ whose canonical key
matches that of $R_v = T-v$. Canonical keys are just serialized coefficient
tensors, so equality of keys means equality in $R_{p,d,n}$. Hence $u = T-v$,
and therefore $u+v=T$.
\end{proof}

\subsection{Multiplicative completion}

The multiplicative completion logic computes the exact quotient $Q_v := T/v$
when it exists, and then fires if either:
\begin{enumerate}[leftmargin=2em]
\item $Q_v$ already appears as an existing node, or
\item $Q_v$ already appears as an existing node.
\end{enumerate}

\begin{proposition}[Multiplicative completion is an exact one-step certificate]
If the multiplicative completion bonus fires, then the exact quotient $Q_v$
already appears as an existing node, and the current state is exactly one
multiplication away from the target.
\end{proposition}

\begin{proof}
By exact quotient soundness, $T = v Q_v$ in $R_{p,d,n}$. If $Q_v$ is already an
existing node, one multiplication of that node with $v$ produces $T$.
\end{proof}

\section{Correctness of Library Accumulation}

After a successful episode, the code registers every non-base node built during
that episode.

\begin{proposition}[Every library entry has a witnessed circuit]
Let $f$ be any polynomial stored in the session library after some number of
episodes. Then there exists a successful episode and a concrete arithmetic
circuit within that episode whose output equals $f$. Moreover, the stored step
count for $f$ is an upper bound on the minimum number of operations needed to
produce $f$ in the environment.
\end{proposition}

\begin{proof}
The library is updated only by \texttt{register\_episode\_nodes()}, which is
called only after \texttt{is\_success=True}. That routine iterates over the
actual node list created during the successful episode and inserts each
non-base node into the library. By the circuit-semantics lemma, each such node
has a witnessed circuit in that episode. The recorded step number is the index
at which the node appeared in that witnessed construction, so it is a valid
upper bound on the optimal construction cost. The library keeps the minimum over
all observed successful episodes, which preserves the upper-bound property.
\end{proof}

\section{Main Corollary}

\begin{corollary}[Soundness summary]
Under the exact semantics of the implemented ring $R_{p,d,n}$:
\begin{enumerate}[leftmargin=2em]
\item the factor library never changes the environment's transition rule or
success predicate;
\item every initial factor subgoal is a genuine nontrivial divisor of the
target;
\item every dynamically added residual or exact quotient subgoal is
algebraically valid;
\item the additive completion bonus is an exact one-step certificate;
\item the multiplicative completion bonus is an exact one-step certificate; and
\item every library entry corresponds to a witnessed reusable subcircuit from a
successful episode.
\end{enumerate}
\end{corollary}

\section{Relation to the Ordinary Polynomial Ring}

The implementation is exact in $R_{p,d,n}$, the truncated quotient ring.
If one wants claims in the ordinary polynomial ring $\Fp[x_0,\dots,x_{n-1}]$
instead, it is enough to assume that all intermediate polynomials appearing in
the trajectories of interest have per-variable degree at most $d$. Under that
assumption, the quotient map $\tau_d$ is injective on those trajectories, so
all equalities proved above in $R_{p,d,n}$ lift to equalities in the usual
polynomial ring.

\section{Practical ICML Framing}

For an ICML submission, a strong theorem statement is therefore:
\begin{quote}
The factor library is a sound reward-guidance mechanism for reusable
subcomputations. It does not alter which action sequences compute the target;
it only adds bonuses attached to algebraically valid factors, residuals, and
exact quotients, together with a session-level memory of witnessed successful
intermediates.
\end{quote}

\paragraph{Scope note.}
This note formalizes the reference factor-library implementation in
\texttt{src/environment/factor\_library.py} and
\texttt{src/environment/circuit\_game.py}. The JAX PPO+MCTS runtime mirrors the
same sound core rewards for initial factor hits, library bonus, additive
completion, multiplicative completion, and direct residual/quotient discovery,
but it does not recursively SymPy-factor every newly discovered
residual/quotient inside the tree-search hot path.

\end{document}
